\clearpage\pagestyle{otMain}\setlength{\otBlocEspacement}{2.5mm}\fontsize{11}{15}\selectfont\corrigesChapitre{20}{Comprendre le hasard}
\begin{blocLivre}{cor-20-1}\begin{corrige}{1}{De l’impossible au certain}\demarche{Compter les faces favorables parmi les 6.}
1. Impossible : $\res0$. 2. Certain : $\res1$. 3. $\res{\frac16}$. 4. $\frac36=\res{\frac12}$.
\end{corrige}\end{blocLivre}
\begin{blocLivre}{cor-20-2}\begin{corrige}{2}{Une urne}\demarche{Issues favorables sur nombre total de boules.}
10 boules en tout. Rouge : $\res{\frac3{10}}$ ; bleue : $\frac5{10}=\res{\frac12}$ ; verte : $\frac2{10}=\res{\frac15}$ ; pas rouge : $\res{\frac7{10}}$.
\label{fig-cor-20-2}\figureGeo[corrige]{hasard-2}
\end{corrige}\end{blocLivre}
\begin{blocLivre}{cor-20-3}\begin{corrige}{3}{Changer d’écriture}\demarche{Passer par les centièmes.}
$\frac14=\res{0{,}25}=\res{25\text{~\%}}$ ; $\frac35=\res{0{,}6}=\res{60\text{~\%}}$ ; $0{,}7=\res{\frac7{10}}=\res{70\text{~\%}}$ ; 15~\% $=\res{\frac{15}{100}}=\res{0{,}15}$.
\end{corrige}\end{blocLivre}
\begin{blocLivre}{cor-20-4}\begin{corrige}{4}{La roue}\demarche{Lister les issues favorables.}
Pair : 2, 4, 6, 8, soit $\frac48=\res{\frac12}$. Au moins 6 : 6, 7, 8, soit $\res{\frac38}$. Multiple de 3 : 3 et 6, soit $\frac28=\res{\frac14}$.
\label{fig-cor-20-4}\figureGeo[corrige]{hasard-4}
\end{corrige}\end{blocLivre}
\begin{blocLivre}{cor-20-5}\begin{corrige}{5}{Une pièce sans mémoire}\demarche{Chaque lancer recommence l’expérience.}
$\res{\text{Non}}$ : la pièce ne se souvient pas des lancers précédents. Au quatrième lancer, la probabilité d’obtenir face vaut toujours $\res{\frac12}$.
\end{corrige}\end{blocLivre}
\begin{blocLivre}{cor-20-6}\begin{corrige}{6}{Deux pièces}\demarche{Compter les chemins de l’arbre.}
4 issues équiprobables : PP, PF, FP, FF. Deux piles : $\res{\frac14}$. Une pile et une face : PF ou FP, soit $\frac24=\res{\frac12}$.
\label{fig-cor-20-6}\figureGeo[corrige]{hasard-6}
\end{corrige}\end{blocLivre}
\begin{blocLivre}{cor-20-7}\begin{corrige}{7}{Deux dés}\demarche{Compter les cases favorables du tableau.}
Somme 7 : six cases, en vert. Probabilité : $\frac6{36}=\res{\frac16}$. Somme 12 : une seule case, 6 et 6 : $\res{\frac1{36}}$.
\label{fig-cor-20-7}\figureGeo[corrige]{hasard-7}
\end{corrige}\end{blocLivre}
\begin{blocLivre}{cor-20-8}\begin{corrige}{8}{Fréquence et probabilité}\demarche{Fréquence observée contre probabilité calculée.}
Fréquence : $\frac{36}{200}=\frac{18}{100}=\res{0{,}18}$. Probabilité : $\frac16\approx\res{0{,}17}$. Les deux nombres sont proches, sans être égaux : c’est normal.
\end{corrige}\end{blocLivre}
\begin{blocLivre}{cor-20-9}\begin{corrige}{9}{Deux urnes}\demarche{Comparer des probabilités, pas des effectifs.}
A : $\frac25=\res{0{,}4}$. B : $\frac38=\res{0{,}375}$. Les chances sont meilleures avec l’urne $\res{\text{A}}$, bien qu’elle contienne moins de boules rouges.
\label{fig-cor-20-9}\figureGeo[corrige]{hasard-9}
\end{corrige}\end{blocLivre}
\begin{blocLivre}{cor-20-10}\begin{corrige}{10}{Une roue inégale}\demarche{Vérifier l’équiprobabilité avant de calculer.}
$\res{\text{Non}}$ : les secteurs ne sont pas de même taille. A occupe la moitié de la roue, d’où une probabilité de $\res{\frac12}$ ; B et C occupent chacun un quart : $\res{\frac14}$ chacun.
\label{fig-cor-20-10}\figureGeo[corrige]{hasard-10}
\end{corrige}\end{blocLivre}
\begin{blocLivre}{cor-20-11}\begin{corrige}{11}{Le jeu de 32 cartes}\demarche{Compter les cartes favorables.}
Cœur : $\frac8{32}=\res{\frac14}$. Roi : $\frac4{32}=\res{\frac18}$. Roi de cœur : $\res{\frac1{32}}$.
\end{corrige}\end{blocLivre}
\begin{blocLivre}{cor-20-12}\begin{corrige}{12}{Le sac de bonbons}\demarche{Appliquer la probabilité au nombre total.}
$0{,}35\times20=\res7$ bonbons à la fraise. Ensuite : 12 bonbons à la fraise sur 25, soit $\frac{12}{25}=\frac{48}{100}=\res{48\text{~\%}}$.
\end{corrige}\end{blocLivre}
\begin{blocLivre}{cor-20-13}\begin{corrige}{13}{La punaise}\demarche{Utiliser la stabilisation des fréquences.}
Fréquence finale : $\frac{620}{1\,000}=\res{0{,}62}$. Comme la fréquence se stabilise, on peut estimer la probabilité à environ $\res{0{,}62}$. On ne peut pas la calculer, car il n’y a pas équiprobabilité.
\label{fig-cor-20-13}\figureGeo[corrige]{hasard-13}
\end{corrige}\end{blocLivre}
\begin{blocLivre}{cor-20-14}\begin{corrige}{14}{Un jeu équitable ?}\demarche{Compter les issues de chaque joueur.}
Alice : 6 issues sur 36. Bob : $1+2+1=4$ issues. Le jeu n’est $\res{\text{pas équitable}}$. En ajoutant la somme 11 (2 issues) à Bob, il a aussi 6 chances sur 36 : le jeu devient équitable.
\end{corrige}\end{blocLivre}
\begin{blocLivre}{cor-20-15}\begin{corrige}{15}{Vrai ou faux}\demarche{Toutes les cases du tableau sont équiprobables.}
$\res{\text{Non}}$ : chacun de ces résultats correspond à une seule case du tableau des 36 issues. Ils ont la même probabilité, $\res{\frac1{36}}$.
\end{corrige}\end{blocLivre}
